\section{Related work}
\label{sec:related}

The closest work to ours is~\cite{Nielsen23cpp},
which proposes a methodology for developing and verifying AMMs in the Coq proof assistant. 
In this approach, an AMM is decomposed into multiple interacting smart contract:
\eg, each minted token is modelled as a single smart contract,
following way fungible tokens are encoded 
in blockchains platforms that do not provide custom tokens natively,
as \eg in Ethereum and Tezos.
These smart contracts are then implemented as Coq functions on top of 
ConCert~\cite{Annenkov20cpp}, a generic model of blockchain platforms and smart contracts mechanized in Coq.
Concert is used to verify behavioural properties of smart contracts,
either in isolation or composed with other smart contracts:
this is fundamental for~\cite{Nielsen23cpp}, 
where AMMs are specified as compositions of multiple contracts. 

More specifically, \cite{Nielsen23cpp} applies the proposed methodology to the Dexter2 protocol, which implements a constant-product AMM based on Uniswap v1
on the Tezos blockchain%
\footnote{\url{https://gitlab.com/dexter2tz/dexter2tz/-/tree/master/}}.
The main properties of AMMs proved in~\cite{Nielsen23cpp} are correspondences
between the state of AMMs and the sequences of transactions executed on the blockchain.
In particular, they prove that:
\begin{itemize}
\item the balance recorded in the main AMM contract is coherent with the actual balance resulting from the execution of the sequence of transactions;
\item the supply of the minted token recorded in the main AMM contract is equal to the actual supply resulting from the execution;
\item the state of the minted token contract is coherent with the execution.
\end{itemize}
Furthermore, starting from the Coq specification of the Dexter2 AMM, \cite{Nielsen23cpp} 
extracts verified CamlLigo code, which is directly deployable on the Tezos blockchain. 

Although both our work and \cite{Nielsen23cpp} involve AMM formalizations within a proof assistant, the ultimate goals are quite different. 
The formalization in~\cite{Nielsen23cpp} closely follows the concrete implementation of a particular AMM instance (Dexter2) and produces a deployable implementation 
that is provably coherent with the proposed Coq specification. 
By contrast, we start from a more abstract specification of AMM,
with the goal of studying the properties that must be satisfied by any implementation coherent with the specification.
%
An advantage of our approach over~\cite{Nielsen23cpp} is that it provides a suitable level of abstraction where proving properties about the economic mechanisms of AMMs,
\ie properties about the gain of users and of equilibria among users' strategies.
In particular, \Cref{th:arbitsol} establishes a paradigmatic property of AMMs, which explains the economic mechanism underlying their design.

Besides these main differences, our AMM model and that in~\cite{Nielsen23cpp}
have several differences.
A notable difference is that our model is based on Uniswap v2, while
\cite{Nielsen23cpp} is based on Uniswap v1.
% \daninote{these main differences ... a first main difference}
In particular, this means that our AMMs can handle arbitrary token pairs, while 
in Dexter2 any AMM pairs a token with the blockchain native crypto-currency. 

At the best of our knowledge, \cite{Nielsen23cpp} and ours are currently 
the only mechanized formalizations of AMMs in a proof assistant.
Many other works study economic properties of AMMs that go beyond those proved in this paper. 
%
The works~\cite{Angeris21analysis} and~\cite{Angeris20aft} study, respectively, an AMM model based on Uniswap similar to ours, and a generalization of the model where the AMM is parameterised over a \emph{trading function} of the AMM reserves, which must remain constant before and after any swap transactions
(in Uniswap v1 and v2, the trading function is just the product between the AMM reserves).
%
The work~\cite{Bartoletti_2022} studies another generalization of Uniswap v2, where the relation between input and output tokens of swap transactions is determined by an arbitrary \emph{swap rate function}, studying the properties of this function that give rise to a sound economic mechanism of AMMs. 
%
While both~\cite{Angeris20aft} and~\cite{Bartoletti_2022} share the common goal
of providing general models of AMMs wherein to study their economic behaviour,
they largely diverge on the formalization:
\cite{Angeris20aft} is based on concepts related to convex optimization problems,
while~\cite{Bartoletti_2022} borrows formalization and reasoning techniques from concurrency theory.
%
The Lean 4 model proposed in this paper follows 
the formalization in~\cite{Bartoletti_2022}, which has the advantage of requiring far less mathematical dependencies: although Mathlib is equipped to reason about convex sets and functions\footnote{\url{https://leanprover-community.github.io/mathlib4_docs/Mathlib/Analysis/Convex/Function}}, it is currently lacking in advanced convex optimization definitions and results as those used in~\cite{Angeris20aft}.

Our AMM model is based on Uniswap v2, one of the most successful AMMs so far.
We briefly discuss some alternative AMM constructions.
%
Balancer~\cite{balancerpaper} generalizes the constant-product function
used by Uniswap to a constant (weighted geometric) mean
$f(r_1,\cdots,r_n) = \prod_{i=1}^n r_i^{w_i}$, 
where the weight $w_i$ reflects the relevance of a token $\tau_i$
in a tuple $(\tau_1,\cdots,\tau_n)$.
% This still fits within the CFMM setting of~\cite{Angeris20aft},
% thus inheriting its results about solvability 
% of the arbitrage problem~\cite{Evans21fc}.
%
Curve~\cite{curvepaper} mixes constant-sum and constant-product functions,
aiming at a swap rate with small fluctuations for large amounts of swapped tokens. 
The work~\cite{Krishnamachari2021fab} studies a variant of the constant-product swap rate invariant, where the rate adjusts dynamically based on oracle prices, 
with the goal of reducing the need for arbitrage transactions.
% but at the cost of lower fee accrual.
%
Other approaches aiming at the same goal are studies in~\cite{virtualAMMbalances,Krishnamachari2021fab}, and implemented in~\cite{moonipaper}.
%
Extending our Lean formalization and results to these alternative AMM designs would require a substantial reworking of our model and proofs.

% bart: rivedere pezzo sotto
% There are several differences between our approach and that 
% of \cite{Danos21defi}, besides the fact that we assume 
% the same swap rate function for all AMMs, and a graph instead of a multi-graph
% (\ie, we admit at most one AMM for each token pair).
% A technical difference is that we assume that the amount $y$ of output 
% tokens received for an amount $x$ of input tokens is given by
% $y = \SX{x,\ammR[0],\ammR[1]} \cdot x$, whereas \cite{Danos21defi}
% defines this amount as $y = f_{\ammR[0],\ammR[1]}(x)$.
% This results in different structural properties for $\SX{x,\ammR[0],\ammR[1]}$
% and $f(x)$ in order to achieve the desired behavioural properties of AMMs.
% Having the AMM reserves $\ammR[0]$, $\ammR[1]$ as parameters of
% our swap rate functions $\SX{}$ has a benefit, in that we can express
% conditions which relate states before and after a transaction:
% this is what happens, \eg, in the additivity, reversibility and homogeneity
% properties (Definitions~\ref{def:sr-additivity},~\ref{def:sr-reversibility}
% and~\ref{def:sr-homogeneity}).
% As a consequence of this choice, compared to~\cite{Danos21defi} 
% our theory encompasses also deposit and redeem actions, 
% providing results that clarify how these actions interfere with swaps
% (\eg, Theorems~\ref{thm:swap-after-dep}, \ref{thm:swap-after-rdm}, \ref{thm:dep-arbitrage}, and \ref{thm:rdm-arbitrage}).
